\section{Conclusion}
\label{sec:conclusion}

To conclude, we have discussed some further subtleties in renormalization
and matching of the quasi-LF correlations on lattice. We proposed a hybrid renormalization procedure to treat the short and long distance correlations separately. The short distance correlations can be renormalized by dividing the same correlator sandwiched in different external states, whereas the long distance ones are renormalized using the Wilson line mass renormalization with a continuity condition to match the short distance region. In this way, we avoid introducing extra non-perturbative effects at large distance in the renormalization stage. We also proposed how to extrapolate to large quasi-LF distance beyond the reach of lattice simulations by utilizing the asymptotic long-range behavior of the correlations, thus avoiding truncations in the ensuing Fourier transform. We finally compared the large-momentum expansion with the CSF approach when applied to LaMET data, showing that the former is a systematic expansion to extract the $x$-dependence of PDFs, whereas the latter is not.
Our proposal here has the potential to greatly improve current computational strategies in lattice applications of LaMET.
